\section{What a Flood CSI Is Actually Worth}
\label{sec:worth}

The most useful thing we can report is not a model score. It is the range a flood score can occupy
on this task \emph{before any model exists}.

\subsection{The score, and the three baselines}
The Critical Success Index is the standard overlap measure in flood mapping. Take the cells the
model says are wet and the cells the radar says are wet. CSI is the number of cells both agree on,
divided by the number of cells either one calls wet. It runs from 0 to 1. Three baselines cost
nothing and bracket it.

\emph{All-wet} predicts every cell is flooded. Its CSI is exactly the observed wet fraction, earned
with zero knowledge. \emph{Random} scatters the correct number of wet cells uniformly at random.
\emph{Climatology} predicts the cells that were wet in most \emph{other} storms of that same area.
It uses no rainfall, no physics and no elevation, and it has one threshold which we fix at 0.5 for
every area rather than tuning it per area, so it is not fitted to the scenes it is scored on. We
also compute a fourth quantity, \emph{SAR versus SAR}, which scores every radar date of an area
against every other radar date of the same area. That is not a baseline. It is an empirical ceiling,
because it measures how much two radar passes of the same city agree with each other.

\begin{table*}[!t]
\centering
\caption{What a CSI is worth on Sentinel-1 flood extent, per area, with no model at all. 150 radar
scenes across 15 areas, of which 149 are scored. Climatology predicts the cells wet in most
\emph{other} storms of that area, with one threshold fixed at 0.5 everywhere. SAR versus SAR scores
every date against every other date of the same area, which is an agreement ceiling between two
radar passes of the same city. Every value is regenerated by \texttt{csi\_net/build\_trivial.py}.}
\label{tab:trivial}
\begin{tabular}{lccccc}
\toprule
area & storms & all-wet & random & climatology & SAR vs SAR \\
\midrule
Patna             &  7 & 0.054 & 0.028 & \textbf{0.286} & 0.214 \\
\midrule
Bengaluru         & 10 & 0.022 & 0.011 & \textbf{0.577} & 0.502 \\
\midrule
Mumbai-South      & 10 & 0.017 & 0.008 & \textbf{0.553} & 0.464 \\
Mumbai-West       & 10 & 0.019 & 0.010 & \textbf{0.464} & 0.387 \\
Mumbai-East       & 10 & 0.020 & 0.010 & \textbf{0.598} & 0.511 \\
Mumbai-North      & 10 & 0.024 & 0.012 & \textbf{0.543} & 0.459 \\
Mumbai-Northeast  & 11 & 0.013 & 0.007 & \textbf{0.391} & 0.271 \\
Mumbai-Harbour    & 11 & 0.021 & 0.010 & \textbf{0.572} & 0.462 \\
\midrule
Chennai           & 10 & 0.021 & 0.011 & \textbf{0.500} & 0.383 \\
\midrule
Doddaballapura    & 10 & 0.014 & 0.007 & \textbf{0.674} & 0.568 \\
Ramanagara        & 10 & 0.009 & 0.004 & \textbf{0.668} & 0.586 \\
Chitradurga       & 10 & 0.019 & 0.010 & \textbf{0.671} & 0.593 \\
Chamarajanagara   & 10 & 0.006 & 0.003 & \textbf{0.522} & 0.464 \\
Kolar             & 10 & 0.030 & 0.015 & \textbf{0.553} & 0.458 \\
Chikkaballapur    & 10 & 0.045 & 0.023 & \textbf{0.604} & 0.523 \\
\midrule
mean              & 149 & 0.022 & 0.011 & \textbf{0.545} & 0.456 \\
\bottomrule
\end{tabular}
\end{table*}

\subsection{The result}
Table~\ref{tab:trivial} is the answer. A climatology with one fixed threshold and no other
parameter scores 0.29 to 0.67, with a mean of 0.545. That is higher than any CSI we are aware of
being reported for a physics-based urban flood model on Sentinel-1, including our own.

It is also higher, \emph{in all fifteen areas}, than the agreement between two Sentinel-1 passes of
that same city. Averaging many passes is more self-consistent than any two of them, which is a
property of the target and not of a model. This caps what any model can honestly be asked to do
here.

The reason the level is so high is that the Sentinel-1 target is dominated by water that has
nothing to do with the storm and that the JRC permanent-water mask does not remove. Tidal flat,
mangrove, seasonal pond and paddy field are all wet on most dates. CSI on this task has therefore
mostly been measuring \emph{can you find standing water}, which our twin never attempted, because
the twin predicts only the water the storm adds.

The signature confirms this reading. On Mumbai-Northeast's one genuine flood, 2026-07-08, with
2\,937 wet cells against a baseline near 450, the climatology scores its \emph{worst} value of
0.075. It wins the quiet dates and loses the flood.

\subsection{What we recommend}
CSI in this literature should always be reported beside all-wet and a leave-one-out climatology
computed on the same scenes. Without them the number is uninterpretable. That applies to the
numbers in the rest of this paper, and in retrospect it applied to the 0.041 we published earlier.
